Machine learning benchmarks drive billions in research investment, yet we lack systematic mechanisms to recognize when a benchmark has exhausted its utility. We demonstrate that benchmark saturation---the point where architectural exploration plateaus---is algorithmically detectable and temporally precedes paradigm shifts by years, not months. Through proof-of-concept experiments on major benchmarks (ImageNet, GLUE, SQuAD), we show that expert consensus on saturation timing exists with >70\% agreement (76-93\% within $\pm$1 year of modal dates), and that simple statistical signals (score convergence via rolling window standard deviation) can detect saturation with statistical significance (Levene's $p<0.05$, 100\% detection rate). Temporal precedence analysis reveals saturations preceded community migration to new paradigms by 2-6 years (mean 48 months, 100\% precedence in tested benchmark-shift pairs), distinguishing internal benchmark exhaustion from external disruption. These findings provide mechanisms for proactive benchmark rotation infrastructure, shifting from reactive organic migration to systematic lifecycle management. Our work operationalizes saturation detection as benchmark governance extension, providing research communities with early warning mechanisms validated on synthetic data (real-world deployment pending Papers With Code data validation).
